\documentclass[10pt,a4paper]{amsart}
\usepackage[margin=19mm]{geometry}
\usepackage{amsmath,amssymb,mathtools}
\usepackage[scr=rsfs,cal=euler]{mathalfa}
\usepackage{xfrac}
\usepackage[unicode,hidelinks,bookmarks=false]{hyperref}
\newcommand{\ie}{i.e.\ }
\newcommand{\cf}{cf.\ }
\newcommand{\resp}{respectively\ }
\newcommand{\Subsection}{\subsection}
\providecommand{\nobreakdash}{\nobreak\hbox{-}}
\setlength{\emergencystretch}{2em}
\multlinegap0pt
\usepackage{bm}
\usepackage{bbm}
\usepackage{dsfont}
\usepackage{xspace}
\usepackage{cancel}
\usepackage{mathtools}
\usepackage{amsthm}
\makeatletter
\@ifundefined{theorem}{\newtheorem{theorem}{Theorem}}{}
\makeatother
\newcommand{\C}{\mathbb C}
\newcommand{\R}{\mathbb R}
\newcommand{\distT}{\operatorname{dist}_{2\pi}}
\title[Counterexamples of Friedlander--Iwaniec dual sums conjecture]{Counterexamples of Friedlander--Iwaniec dual sums conjecture}
\author{Khai-Hoan Nguyen-Dang}
\date{Source: arXiv:2607.16695v1 (CC BY 4.0)}
\begin{document}
\maketitle

\noindent\emph{Source and licence.} Counterexamples of Friedlander--Iwaniec dual sums conjecture. Version arXiv:2607.16695v1 (CC BY 4.0), distributed under the Creative Commons Attribution 4.0 International licence, as recorded in the arXiv licence metadata for this exact version and shown on the abstract page https://arxiv.org/abs/2607.16695v1. Full rights notice: licenses/arxiv-2607.16695-CC-BY-4.0.txt.\par\smallskip

\noindent\textit{Editorial scope.} Counted unit: the Theorem whose source label is \texttt{thm:sector} in the article (source0.tex lines 587--620, source label \texttt{thm:sector}), delivered together with the article's own complete original proof of it (source0.tex lines 622--695, one \texttt{proof} environment of 74 lines). No mathematics is written, repaired or re-derived: statement and proof are the article's own text line for line. Required context retained uncounted: none. Every definition, display and label of those lines is retained unchanged. Omitted from this package and not redistributed: 1--586, 621--621, 696--1181. Nothing inside a retained excerpt is omitted or altered: every retained body is delivered exactly as the article's lines are written, so no omission receipt of any kind accompanies this package. Citations: the retained excerpts carry no citation command at all, so no bibliography entry of the article is transcribed and no reference is redistributed. Reference adaptation: none was needed and none was made; every reference inside the delivered unit resolves inside it, so no unresolved reference or citation is produced. Printed slips kept literal: no printed word, symbol, hypothesis or number of the counted statement or of its proof was altered. Layout and class: the article's own class is \texttt{amsart} with packages including fontenc, lmodern, microtype, amsmath, amssymb, amsthm, mathtools, mathrsfs, enumitem, hyperref; the wrapper is the lane's pinned \texttt{amsart} editorial wrapper at 10pt on A4, which restates the theorem-like environments the retained text needs and adds the article's own macro definitions that the retained text needs (\texttt{\textbackslash C}, \texttt{\textbackslash R}, \texttt{\textbackslash distT}), with extra wrapper packages (bm, bbm, dsfont, xspace, cancel, mathtools). The delivered numbering is the unit's own. Assets: no retained span contains a figure environment, a graphics-inclusion command, a PDF-page inclusion, a table or a TikZ picture; no image file is copied into this package, the package declares no asset dependency and no image file is redistributed with it. Licence: version arXiv:2607.16695v1 of this article is distributed under the Creative Commons Attribution 4.0 International licence, as recorded in the arXiv licence metadata for this exact version and shown on the abstract page https://arxiv.org/abs/2607.16695v1. A full rights notice is delivered at licenses/arxiv-2607.16695-CC-BY-4.0.txt. Characters: every retained excerpt is pure ASCII, so the delivered engine, XeLaTeX with its default Latin Modern text font, reports no missing character.
\begin{theorem}\label{thm:sector}
Let $0<\nu<1$, let $\beta\in\R$, and let $\kappa>0$.  Let $\Omega:\R\to\C$ be a continuous, $2\pi$-periodic function that is not identically zero, and let $(c_n)_{n\ge1}$ be real numbers satisfying
\begin{equation}\label{eq:coeff-lower}
  c_n\ge c_*>0\qquad(n\ge1).
\end{equation}
For $x>0$ and $M\ge0$, define
\begin{equation}\label{eq:general-S}
  S_{\Omega}(x,M)
  :=\sum_{n\le M}c_n n^{-\beta}
  \Omega\!\left(\kappa(nx)^\nu\right).
\end{equation}
Fix $\lambda>0$ and
\begin{equation}\label{eq:a-general-range}
  -1<a<\nu^{-1}-1.
\end{equation}
Then, for every sufficiently large integer $N$, there exist $x_N>0$ and an integer $H_N$ with $1\le H_N<N/2$ such that
\begin{align}
  x_N
  &=\lambda N^a\left(1+O\bigl(N^{-\nu(a+1)}\bigr)\right),\label{eq:sector-x}\\
  H_N
  &\asymp N^{1-\nu(a+1)},\label{eq:sector-H}
\end{align}
and
\begin{equation}\label{eq:sector-block}
  \bigl|S_{\Omega}(x_N,N)-S_{\Omega}(x_N,N-H_N)\bigr|
  \gg N^{1-\beta-\nu(a+1)}.
\end{equation}
Consequently, for some $M_N\in\{N-H_N,N\}$,
\begin{equation}\label{eq:sector-max}
  |S_{\Omega}(x_N,M_N)|
  \gg N^{1-\beta-\nu(a+1)}.
\end{equation}
All implied constants in this theorem may depend on the fixed data $\nu,\beta,\kappa,\lambda,a,\Omega$, and $c_*$.
\end{theorem}

\begin{proof}
Choose $t_0\in[0,2\pi)$ with $\Omega(t_0)\ne0$, and set
\[
  \xi=\frac{\overline{\Omega(t_0)}}{|\Omega(t_0)|}.
\]
By continuity on the circle, there is a number $\delta\in(0,\pi)$ such that
\begin{equation}\label{eq:sector-arc}
  \operatorname{Re}\bigl(\xi\Omega(t)\bigr)
  \ge\frac12|\Omega(t_0)|
  \qquad
  \text{if }\distT(t,t_0)\le\delta.
\end{equation}
Put
\[
  Y_N=\kappa\lambda^\nu N^{\nu(a+1)}
\]
and choose
\begin{equation}\label{eq:j-choice}
  j_N=\left\lceil\frac{Y_N-t_0}{2\pi}\right\rceil,
  \qquad
  x_N=\frac1N
  \left(\frac{2\pi j_N+t_0}{\kappa}\right)^{1/\nu}.
\end{equation}
Since $2\pi j_N+t_0=Y_N+O(1)$ and $Y_N\to\infty$, the mean-value theorem applied to $u\mapsto u^{1/\nu}$ gives \eqref{eq:sector-x}.  Moreover,
\begin{equation}\label{eq:exact-resonance}
  \kappa(Nx_N)^\nu=2\pi j_N+t_0
  \equiv t_0\pmod{2\pi}.
\end{equation}

Let
\[
  \Phi_N(u)=\kappa(ux_N)^\nu.
\]
For $u\in[N/2,N]$, equation \eqref{eq:sector-x} gives
\begin{equation}\label{eq:general-derivative}
  0<\Phi_N'(u)
  =\kappa\nu x_N^\nu u^{\nu-1}
  \le C N^{\nu(a+1)-1}
\end{equation}
for all sufficiently large $N$, where $C$ depends only on the fixed data.  Choose $\eta>0$ so small that $2C\eta\le\delta$, and set
\begin{equation}\label{eq:H-choice-general}
  H_N=\left\lfloor\eta N^{1-\nu(a+1)}\right\rfloor.
\end{equation}
The two inequalities in \eqref{eq:a-general-range} give
$0<1-\nu(a+1)<1$.  Hence $H_N\to\infty$ and $H_N=o(N)$, so $1\le H_N<N/2$ for all sufficiently large $N$.  If $N-H_N<n\le N$, then by \eqref{eq:general-derivative},
\[
  |\Phi_N(N)-\Phi_N(n)|
  \le C H_NN^{\nu(a+1)-1}
  \le\delta.
\]
Together with \eqref{eq:exact-resonance} and \eqref{eq:sector-arc}, this yields
\begin{equation}\label{eq:sector-positive}
  \operatorname{Re}\bigl(\xi\Omega(\Phi_N(n))\bigr)
  \ge\frac12|\Omega(t_0)|
  \qquad(N-H_N<n\le N).
\end{equation}
On this block, $n\asymp N$, and therefore $n^{-\beta}\asymp_{\beta}N^{-\beta}$.  Using \eqref{eq:coeff-lower} and \eqref{eq:sector-positive}, we obtain
\begin{align*}
  &\bigl|S_{\Omega}(x_N,N)-S_{\Omega}(x_N,N-H_N)\bigr|\\
  &\quad\ge
  \operatorname{Re}\!\left(
    \xi\sum_{N-H_N<n\le N}
    c_n n^{-\beta}\Omega(\Phi_N(n))
  \right)\\
  &\quad\gg H_NN^{-\beta}
  \gg N^{1-\beta-\nu(a+1)}.
\end{align*}
This proves \eqref{eq:sector-block}.  Finally,
\[
  |S_{\Omega}(x_N,N)-S_{\Omega}(x_N,N-H_N)|
  \le |S_{\Omega}(x_N,N)|+|S_{\Omega}(x_N,N-H_N)|,
\]
which implies \eqref{eq:sector-max}.
\end{proof}

\end{document}
